\section{(H) FUNCTIONS}

PROPOSITION 3. If \(\mathfrak{K}_0\) is a Hardy field there exists a Hardy field \(\mathfrak{K}\) containing \(\mathfrak{K}_0\) and such that, for every function \(z \in \mathfrak{K}\) , not identically zero on a neighbourhood of \(+\infty\) , both \(e^z\) and \(\log |z|\) belong to \(\mathfrak{K}\) .

Denote by \(\mathfrak{K}\) the set of functions \(f\in \mathcal{H}(\mathfrak{F},\mathbf{R})\) having the following properties: for each function \(f\in \mathfrak{K}\) there is a finite number of Hardy fields \(\mathfrak{K}_1,\mathfrak{K}_2,\ldots ,\mathfrak{K}_n\) (the number \(n\) and the fields \(\mathfrak{K}_i\) depending on \(f\) ) such that \(f\in \mathfrak{K}_n\) and, for \(0\leqslant i\leqslant n - 1\) , one has \(\mathfrak{K}_{i + 1} = \mathfrak{K}_i(u_{i + 1})\) where \(u_{i + 1}\) is equal either to \(e^{z_i}\) or to \(\log |z_i|\) with \(z_{i}\) belonging to \(\mathfrak{K}_i\) and not vanishing identically on a neighbourhood of \(+\infty\) . One says that \(u_{1},u_{2},\ldots ,u_{n}\) form a definition sequence for the field \(\mathfrak{K}_n\) and of the function \(f\) ; the same function can naturally admit several definition sequences.

By def. 1 of V, p.~247, every function \(f \in \mathfrak{K}\) , not identically zero on a neighbourhood of \(+\infty\) , has constant sign and is differentiable on an interval \([x_0, +\infty[\) ; if \(f \in \mathfrak{K}_n\) , then \(1/f\) and \(f'\) are equal to functions in \(\mathfrak{K}_n\) , thus to functions in \(\mathfrak{K}\) , on a neighbourhood of \(+\infty\) . To see that \(\mathfrak{K}\) is a Hardy field it is enough to prove that if \(f\) and \(g\) are two functions in \(\mathfrak{K}\) then \(f - g\) and \(fg\) are equal to functions in \(\mathfrak{K}\) on a neighbourhood of \(+\infty\) . Now let \(u_1, u_2, \ldots, u_m\) be a definition sequence for \(f\) , and \(v_1, v_2, \ldots, v_n\) a definition sequence for \(g\) . The sequence \(u_1, u_2, \ldots, u_m, v_1, v_2, \ldots, v_n\) obtained by concatenating the sequences \((u_i)\) and \((v_i)\) is again a definition sequence of a Hardy field \(\mathfrak{K}_{m+n}\) and this field contains \(f\) and \(g\) , so \(f - g\) and \(fg\) are equal to functions in \(\mathfrak{K}_{m+n}\) on a neighbourhood of \(+\infty\) .

One says that the Hardy field \(\mathfrak{K}\) defined in the proof of prop. 3 is the (H) extension of the Hardy field \(\mathfrak{K}_0\) .

If \(\mathfrak{K}'\) is another Hardy field possessing the properties stated in prop. 3, it follows from the construction of \(\mathfrak{K}\) that \(\mathfrak{K} / \mathrm{R}_{\infty}\) is contained in \(\mathfrak{K}' / \mathrm{R}_{\infty}\) . By abuse of language, one says that the (H) extension of the Hardy field \(\mathfrak{K}_0\) is the smallest Hardy field \(\mathfrak{K}\) having these properties.

DEFINITION 2. The field of (H) functions is called the (H) extension of the Hardy field \(\mathbf{R}(x)\) of rational functions with real coefficients. Any function belonging to this extension is called an (H) function.

Following this definition, if f is an (H) function, not identically zero on a neighbourhood of \(+\infty\) , then \(e^{f}\) and \(\log |f|\) are also (H) functions. More generally, if g is a second (H) function, and \(u_{1}, u_{2}, \ldots, u_{n}\) a definition sequence for g, and if \(f(x)\) tends to \(+\infty\) with x, one sees by induction on n that the composite functions \(u_{1} \circ f, u_{2} \circ f, \ldots, u_{n} \circ f\) and \(g \circ f\) are (H) functions.

